\section{Exact-Action Compatibility and New-Context Transfer}
\label{sec:pending-studies}
The follow-up connects the class-risk theorem to the retained representations and evaluates transport under a new language construction. Its five analyses completed without dropping declared conditions: 48 spectral cells, 48 constructed witnesses, all 2,952 image-subgroup units, 36 fixed mechanism settings, and 48 new-context cells. The methods and provenance are in \cref{app:pending-protocols}; the results below concern distinct targets rather than a pooled success score.

\subsection{A feasible exact action need not be the best predictor}
\input{generated/followup_spectral_bounded.tex}
The empirical spectral budget averages $0.98510$ for weekdays and $0.93528$ for corners. The corresponding mean cell-wise lower/upper bounds are $[0.02980,0.05863]$ and $[0.12944,0.24895]$ in the test bank's own second-moment-whitened metric (\cref{tab:followup-spectral}). These are deterministic bounds on the specified finite-bank action class. They establish compatibility at those risk levels, not existence of a native semantic computation.

Trivial-type components account for selected energy $b_{\rm triv}=0.61540/0.72742$ on average. Thus much of the retained budget is state-invariant, although a trivial component need not be purely contextual and a nontrivial component need not reflect computation. The real-type correction is at numerical scale in 47/48 conditions (absolute correction at most $5.33\times10^{-15}$). The exception is Llama-3.1-8B, C7, layer 13: $\Delta_{\rm type}=0.13235$, or $\Delta_{\rm type}/16=0.00827$. The typed restriction is essential to the theorem and counterexample, but changes this rank-16 panel materially in only one condition.

\input{generated/followup_witness_bounded.tex}
All 48 constructed witnesses satisfy the frozen action-law and orthogonality checks, and their construction-bank risks lie within the analytic bounds. Frozen test losses average $0.02806/0.11117$, compared with $0.02141/0.08878$ for independent OLS (\cref{tab:followup-witness}). The witness is worse than OLS and fixed-coefficient ridge in every condition. This is an attained-risk comparison, not a measurement of the optimal risk premium: the witness need not minimize $L^\star$. These prediction losses use \emph{training} whitening for both construction and test banks, whereas \cref{tab:followup-spectral} uses the test bank's own metric. The metrics also define different admissible classes of raw-coordinate functions. Conjugating the frozen training action into test-whitened coordinates preserves its group law but need not preserve orthogonality (\cref{app:metric-conjugation}); its transfer loss is therefore not checked against the test-bank exact-orthogonal bounds.

\begin{table}[t]
\centering\small
\input{generated/bounded_risk_reference_table.tex}
\caption{Same-metric reference and class-risk premium on the original Study~2 test bank. Each domain averages \BoundedCellsPerDomain{} fixed cells equally. $p^-=\max\{0,2(1-\gamma_r)-\eta\}$ and $p^+=2(1-\gamma_r^2)-\eta$ are computed per cell before averaging. They bound the mean cost of imposing exact laws and orthogonality jointly. $L(I)$ uses total trivial energy and was also checked by direct residuals on the saved coordinates. These are finite-bank quantities, not transfer losses or confidence limits.}
\label{tab:bounded-risk}
\end{table}

\paragraph{The class constraint has a measurable cost on these banks.}
\Cref{tab:bounded-risk} reports the same-metric identity reference and the bounds on $L^\star-\eta$ from \cref{eq:risk-premium-main}. The mean premium lies in $[\BoundedPremiumLowerCseven,\BoundedPremiumUpperCseven]$ for C7 and $[\BoundedPremiumLowerDfour,\BoundedPremiumUpperDfour]$ for D4. All \BoundedAllPositive{} lower endpoints exceed the original numerical tolerance. This fixed-panel count rules out matching independent linear prediction within the specified exact orthogonal class at those bounds; it is not a significance test or a conclusion about unseen contexts. Identity risks average $\BoundedIdentityCseven/\BoundedIdentityDfour$, illustrating the gain available from state-dependent prediction even though much total energy is invariant. The calculation constrains the whole class and is separate from the constructed witness's higher test loss. Complete cell and type values are in \cref{app:bounded-risk-details}.

\subsection{High compatibility does not uniquely favor the specified image}
The enumeration evaluates 120 distinct cyclic order-seven permutation images for C7 and three dihedral order-eight images for D4, on each of 24 model/readout banks. The specified semantic image has uniquely largest $\gamma_r$ in 7/24 and 16/24 conditions, respectively. C7's full-image gamma spread has median $3.88\times10^{-5}$ and maximum $4.07\times10^{-4}$; for D4 these are $1.36\times10^{-4}$ and $0.03984$. In C7, ranks can therefore differ substantially while the budget changes little. The specified C7 image ranges from rank 1 to a tie interval of 102--105 out of 120; D4 includes all three ranks.

\begin{table}[t]
\centering\small
\input{generated/bounded_registration_margin_table.tex}
\caption{Absolute spectral-budget margins of the specified image against the strongest and median alternatives, excluding that image itself. Summaries use all \BoundedCellsPerDomain{} cells per domain. U/T/B counts unique leading, tied leading and behind under the original pairwise tolerance. Raw margins retain their signs; they are not permutation $p$-values, exact-action risk rankings or probabilities of semantic correctness.}
\label{tab:bounded-margins}
\end{table}
For semantic image $H_0$, define $d_{\rm top}=\gamma(H_0)-\max_{H\ne H_0}\gamma(H)$ and $d_{\rm median}=\gamma(H_0)-\operatorname{median}_{H\ne H_0}\gamma(H)$. \Cref{tab:bounded-margins} adds their absolute scale to the ranks. C7's median margin is negative against the best alternative but positive against the median alternative; its small margins show why ranking alone can overstate separation. D4's positive median coexists with both larger gains and losses across cells. These observations do not establish statistical equivalence or a uniformly preferred semantic image. Full margins and unchanged tolerance rules are in \cref{app:bounded-margin-details}.

These counts are descriptive across selected conditions. The comparison fixes coordinates, metric, weights and complete source/action coverage, but each image defines different target correspondences. It ranks spectral compatibility with registered tasks, not the exact optimum $L^\star$, a shared set of target labels, or semantic posterior probability. Even a uniquely leading image leaves action automorphisms and state-code mechanisms possible. The mechanism controls make the latter concrete: both the noiseless shared-action mechanism M1 and stable-code mechanism M3 yield machine-scale witness loss; the sampled contextual-orientation mechanism M2 has appreciably lower gamma (\cref{tab:followup-simulations}). These 36 fixed settings illustrate the constructions rather than calibrate a semantic classifier.

\subsection{The transport advantage changes under a new construction}
\input{generated/followup_context.tex}
The new bank has 64 assistant-authored paired frames per domain, producing 896 C7 and 512 D4 prompts per model. Both primary predictors are fitted on old training data and selected using the same old-validation context-difference loss. New prompts are used for neither fitting nor selection. The earlier coordinate comparator was selected on coordinate loss; the retained secondary arm using that rule also shows the same domain signs here. Thus the comparison with the old study is not a controlled language-only intervention. They change the outer narrative and limited fixed factual wording together; the state vocabulary is unchanged. This corrects the frozen protocol's overly narrow description of an outer-narrative-only change.

The primary equal-domain contrast, difference fitting minus coordinate fitting, is $+0.000178$, with recorded family-resampling percentiles $[-0.001035,+0.001362]$. The domain means point in opposite directions: $-0.005833$ for C7 and $+0.006190$ for D4. Difference fitting improves 20/24 weekday conditions but only 8/24 corner conditions. Identity has the lowest mean C7 loss, while both fits improve on identity for D4. The nearly cancelling pooled contrast establishes neither equivalence nor a general advantage for difference fitting.

\begin{figure}[t]
\centering
\includegraphics[width=.88\textwidth]{figures/context_transfer_followup.pdf}
\caption{New-context contrast by model and domain, averaging four fixed readouts. Negative values favor difference fitting. All six C7 model means are negative; D4 reverses direction in four models. Points are descriptive model means, not independent replications or confidence intervals. The primary aggregate and paired-family sensitivity ranges appear in \cref{tab:followup-context}.}
\label{fig:context-followup}
\end{figure}
The result narrows the original construction's finding: transport depends on the fitting objective, domain and language construction. The bank tests neither unseen state vocabulary nor arbitrary human prose; post-target text cannot affect a causal decoder's state-token readout. The family bootstrap preserves repeated cells and domains together and describes this authored bank, as specified in \cref{sec:setup}. No new significance, equivalence or native causal conclusion is drawn.
